\chapter{Stand der Technik}\label{stand-der-technik}

Dieses Kapitel arbeitet die technologischen Grundlagen auf, die das in Kapitel~\ref{motivation-und-anforderungen} beschriebene System erfordert. Es behandelt Large Language Models (Abschnitt~\ref{large-language-models}), Prompt Engineering (Abschnitt~\ref{prompt-engineering}), automatisierte Sportberichterstattung (Abschnitt~\ref{automatisierte-sportberichterstattung}), ETL-Workflows (Abschnitt~\ref{etl-prozesse-und-workflow-automatisierung}) und den Webanwendungs- und Infrastruktur-Stack (Abschnitt~\ref{backend-frontend-und-deployment-technologien}).

\section{Large Language Models}\label{large-language-models}

Large Language Models (LLMs, dt. Große Sprachmodelle) werden auf umfangreichen Textkorpora trainiert und können natürliche Sprache verstehen und generieren. Technisch basieren sie auf der \textbf{Transformer-Architektur} (Vaswani et al.~2017), die Text mithilfe von \textbf{Aufmerksamkeitsmechanismen} (Attention) verarbeitet: Das Modell gewichtet bei der Verarbeitung eines Tokens alle anderen Token im Kontext nach ihrer Relevanz, wodurch es semantische Zusammenhänge über lange Sequenzen erfasst. Text wird autoregressiv generiert. Jedes neue Token entsteht auf Basis des bereits generierten Kontexts. LLMs unterstützen zudem \textbf{In-Context Learning}: Modelle lassen sich über Beispiele im Prompt konditionieren, ohne ihre Gewichte anzupassen (vgl. Brown et al.~2020, S.~4).

LLMs unterscheiden sich zudem erheblich in Größe und Inferenzaufwand. \textbf{Kompakte Modelle} können bei strukturierten, auf kurze Ausgaben ausgelegten Aufgaben eine vergleichbare Ausgabequalität bei deutlich niedrigerer Latenz und geringeren Kosten als große Modelle erreichen. Dieser Zusammenhang wird in aktuellen anbieterseitigen Evaluations- und Benchmark-Berichten beschrieben (vgl. OpenAI 2023; Mistral AI 2024; Google 2024). Die Modellwahl ist damit eine eigenständige Designentscheidung mit direktem Einfluss auf Latenz- und Kostenbudget.

Für den vorliegenden Anwendungsfall sind \textbf{LLM-API-Aggregatoren} als Infrastrukturkategorie relevant: Dienste wie OpenRouter stellen eine einheitliche API-Schnittstelle bereit, die Modelle mehrerer Anbieter hinter einem einzigen Endpunkt bündelt. Ein Modellwechsel erfordert damit lediglich eine Konfigurationsänderung statt einer Anpassung der Anwendungslogik. Diese Abstraktion entkoppelt die Anwendung von einzelnen Anbietern, ermöglicht Fallback-Mechanismen bei Verfügbarkeitsausfällen und vereinfacht Modellvergleiche auf identischer Prompt-Basis. Die architektonische Einbettung dieses Ansatzes wird in Abschnitt~\ref{multi-provider-architektur} konkretisiert.

Trotz ihrer Leistungsfähigkeit weisen LLMs strukturelle Schwächen auf, insbesondere die Neigung zu \textbf{Halluzinationen}, also zur Generierung sachlich unzutreffender oder erfundener Inhalte (OpenAI 2023, S.~46). Das Fehlerspektrum reicht dabei von frei erfundenen Sachverhalten (faktische Halluzinationen) bis zur fehlerhaften Interpretation vorhandener Eingabedaten (kontextuale Halluzinationen). Besonders kritisch sind frei erfundene Aussagen, weil sie ohne externen Faktenabgleich schwer zu erkennen sind und im journalistischen Kontext die publizistische Glaubwürdigkeit gefährden (vgl. Thürbach, zit. n. Bluhm \& Schäfer 2023, S.~33). Gegenmaßnahmen umfassen redaktionelle Kontrollinstanzen und explizite Schutzregeln in der Prompt-Gestaltung (vgl. Abschnitt~\ref{prompt-engineering-strategie}).





\section{Prompt Engineering}\label{prompt-engineering}

Die Qualität der Ausgaben eines LLM hängt maßgeblich vom \textbf{Prompt}
ab. Prompt Engineering bezeichnet die systematische Gestaltung dieser
Eingaben, um das Modell ohne Änderung seiner Gewichte zu einem
gewünschten Verhalten zu lenken (Brown et al.~2020).

Für die Anwendung von LLMs auf spezifische Aufgaben existieren grundsätzlich zwei Ansätze: \textbf{Fine-tuning} erfordert ein aufgabenspezifisches Nachtrainieren des Modells auf kuratierten Datensätzen und ist ressourcenintensiv (Rechenzeit, Infrastruktur, Datenkuration). \textbf{Prompting-basierte Ansätze} hingegen nutzen das vortrainierte Modell direkt und steuern dessen Verhalten ausschließlich über die Eingabeformulierung (vgl. Brown et al.~2020, S.~3 u.~6). Für Domänen mit fehlenden oder unzureichenden Trainingsdaten ist Prompting in der Regel methodisch vorzuziehen: Fine-tuning ist dann häufig mit erhöhtem Risiko eingeschränkter Generalisierbarkeit verbunden, während Prompting-basierte Ansätze mit wenigen kuratierten Beispielen auskommen und schneller iterierbar bleiben (vgl. Goodfellow et al.~2016, S.~111--117).

Innerhalb des Prompting-Paradigmas lassen sich drei wesentliche Techniken unterscheiden. Beim \textbf{Zero-Shot Prompting} erhält das Modell ausschließlich eine Aufgabenbeschreibung ohne Demonstrationsbeispiele (vgl. Brown et al.~2020, S.~7). Beim \textbf{Few-Shot Prompting} werden zusätzlich wenige Beispiele als Stilreferenz bereitgestellt (vgl. Brown et al.~2020, S.~6), die das Modell auf eine gewünschte Tonalität konditionieren, ohne dass Fine-tuning erforderlich ist. Beim \textbf{Role Prompting} wird dem Modell eine explizite Persona zugewiesen (z.\,B. „Du bist ein Sportjournalist''), die es auf eine konsistente sprachliche Identität verpflichtet. Rollenbasierte Instruktionen verbessern nachweislich die Ausgabekonsistenz bei stilistisch definierten Aufgaben (vgl. Shanahan et al.~2023, S.~3). Instruction-optimierte Modelle folgen solchen Anweisungen zuverlässig auch ohne aufwändige Prompt-Konstruktion (vgl. Ouyang et al.~2022, S.~3).

Moderne Chat-Modelle unterscheiden strukturell zwischen zwei
Prompt-Rollen (OpenAI 2024, Anthropic 2024). Der \textbf{System Prompt} legt die übergeordnete Persona
und Constraints des Modells fest (z.\,B. Stilanforderungen, inhaltliche
Grenzen, Ausgabeformat), während der \textbf{User Prompt} die konkrete
Aufgabe mit aufgabenspezifischen Daten enthält. Diese Trennung erlaubt
es, allgemeine Verhaltensregeln einmalig im System Prompt zu definieren
und den User Prompt auf die jeweilige Aufgabe zu beschränken.

Ein weiterer zentraler Steuerungsparameter ist die \textbf{Temperature},
die die Wahrscheinlichkeitsverteilung der Token-Auswahl beeinflusst
(vgl. Brown et al.~2020). Niedrige Werte (z.\,B. 0{,}0--0{,}3) führen zu deterministischeren, faktentreueren Ausgaben. Höhere Werte (z.\,B. 0{,}7--1{,}0) erhöhen die kreative Varianz. Je nach Aufgabentyp ist daher ein unterschiedlicher Temperaturbereich geeignet: Generierungsaufgaben mit Stilanforderungen profitieren von mittleren Werten, Übersetzungsaufgaben mit strengen Treuekriterien von niedrigen. Die konkrete Parameterwahl des vorliegenden Systems wird in Abschnitt~\ref{prompt-engineering-strategie} begründet.



\section{Automatisierte Sportberichterstattung}\label{automatisierte-sportberichterstattung}

Die automatisierte Erzeugung natürlichsprachlicher Texte aus strukturierten Daten (Data-to-Text Generation) bildet die technische Grundlage für KI-gestützte Berichterstattungssysteme (vgl. Puduppully \& Lapata 2021, S.~510). Die Systementwicklung lässt sich in drei Generationen beschreiben. Regelbasierte Template-Architekturen lieferten vorhersagbare, aber stilistisch monotone Ausgaben. Neuronale End-to-End-Systeme überwanden die Monotonie durch datengetriebenes Lernen, erforderten jedoch umfangreiches domänenspezifisches Training und tendierten zu faktischen Inkonsistenzen (vgl. Puduppully \& Lapata 2021, S.~520, Boisserée, zit. n. Beils 2023, S.~207). LLM-basierte Systeme ermöglichen als dritte Generation eine flexible Stilkonditionierung durch Prompt-Design ohne domänenspezifisches Training und machen damit den stilistischen Vorteil neuronaler Systeme ohne deren Trainingsaufwand zugänglich. Das Halluzinationsrisiko bleibt als strukturelles Restproblem bestehen (vgl. Abschnitt~\ref{large-language-models}).

Trotz dieser Entwicklung sind in der redaktionellen Praxis weiterhin templatebasierte Ansätze verbreitet: \textit{Die Welt} etwa erzeugt automatisierte Spielberichte via Retresco (vgl. Beils 2023, S.~207), weil Zuverlässigkeit dort höher gewichtet wird als stilistische Varianz.

Empirische Befunde zur automatisierten Sportberichterstattung zeigen, dass Leser KI-generierte Sporttexte in Lesbarkeit und Vertrauenswürdigkeit ähnlich wie menschlich verfasste Texte bewerten, diese jedoch in wahrgenommener Expertise und Stilqualität schlechter einschätzen (Clerwall 2014, S.~529). Der Befund ist vor dem Hintergrund zu interpretieren, dass Clerwall eine frühe, vor-LLM-basierte Technologiegeneration untersucht; die Modellqualität und stilistische Varianz haben sich seither deutlich erhöht. Graefe et al.\ (2018, S.~7) replizieren den Kernbefund für strukturierte Nachrichtendomänen und zeigen, dass die wahrgenommene Qualitätslücke beim faktenintensiven Sportbericht schmaler ausfällt als bei allgemeinen Nachrichtentexten. Dies deutet darauf hin, dass das strukturierte, datendichte Sportformat algorithmischer Generierung besonders entgegenkommt. Für das Genre Liveticker, das von sprachlicher Lebendigkeit, Ellipsen und situativer Einordnung lebt, ist dieser Stilnachteil besonders relevant: Formelhaft wirkende, templatebasierte Texte verstoßen gegen Genrekonventionen und senken die Leserbindung.

Die Befunde zeigen zugleich, in welchen Bereichen menschliches Urteil weiterhin zentral bleibt. Das Konzept des \textbf{Human-in-the-Loop (HITL)} bezeichnet Systemdesigns, in denen menschliche Akteure gezielt in automatisierte Entscheidungsprozesse eingebunden bleiben (Mosqueira-Rey et al.~2023, S.~2). Für den journalistischen Kontext ist es besonders relevant, weil die publizistische Verantwortung — die Festlegung von Richtigkeit, Ton und Relevanz — in diesem Anwendungskontext nicht vollständig an ein Modell delegiert werden kann, das weder redaktionellen Kontext noch die Implikationen einer Falschaussage verlässlich einschätzen kann (vgl. Canalp, zit. n. Beils 2023, S.~57). Die systemseitige Umsetzung dieses Prinzips ist in Abschnitt~\ref{status-lifecycle-und-betriebsmodi} beschrieben.

Das \textbf{LLM-as-Judge-Verfahren} (Zheng et al.~2023) setzt ein vom generierenden Modell verschiedenes Sprachmodell als automatisierten Bewerter ein: Es beurteilt KI-generierte Ausgaben entlang definierter Qualitätsdimensionen, ohne dass menschliche Annotationen für jeden Bewertungsschritt erforderlich sind. Korrelationsstudien zeigen eine moderate Übereinstimmung mit menschlichen Urteilen bei deutlich geringerem Annotationsaufwand (vgl. Zheng et al.~2023, S.~8). Im Kontext automatisierter Sportberichterstattung eignet sich das Verfahren als skalierbares Evaluationsprotokoll für Dimensionen wie Korrektheit, Tonalität und Vollständigkeit. Zugleich sind mögliche Bewerterverzerrungen zu berücksichtigen, deren Konsequenzen in Abschnitt~\ref{methodische-einordnung} diskutiert werden. Die Anwendung im vorliegenden Evaluationsdesign wird in Abschnitt~\ref{erweiterte-evaluation} dargestellt.






\section{ETL-Prozesse und Workflow-Automatisierung}\label{etl-prozesse-und-workflow-automatisierung}

Die Datenbeschaffung in datengetriebenen Anwendungen folgt überwiegend dem \textbf{Extract-Transform-Load-Prinzip (ETL)}: Daten werden aus heterogenen Quellen extrahiert, in ein einheitliches Zielformat transformiert und in den Datenspeicher geladen (vgl. Freitas et al.~2025, S.~807).

Workflow-Orchestrierungsplattformen modellieren ETL-Logik als konfigurierbare DAG-Prozesse, die Abhängigkeiten zwischen Verarbeitungsschritten explizit machen (vgl. Venkiteela 2025, S.~3). Low-Code-Varianten dieser Plattformen reduzieren die Einstiegshürde gegenüber prozeduralem ETL-Code und erlauben bei Bedarf die Einbettung von Custom-Code-Blöcken (vgl. Venkiteela 2025, S.~2).

Innerhalb dieser Klasse lassen sich cloud-basierte iPaaS-Dienste (z.\,B. Zapier, Make) von self-hostbaren Open-Source-Plattformen (z.\,B. n8n) unterscheiden (vgl. Venkiteela 2025, S.~2). Cloud-Dienste bieten sofortige Verfügbarkeit ohne Infrastrukturaufwand, unterliegen jedoch Datenschutzeinschränkungen und eingeschränkter Kontrolle über Ausführungsumgebung und Kosten. Self-hostbare Plattformen ermöglichen dagegen den Betrieb in der eigenen Infrastruktur mit voller Kontrolle über Datenflüsse (vgl. Nakavisute \& Sincharoonsak 2025).

Für Echtzeitanwendungen wie Liveticker ist die Latenz der Orchestrierung ein zusätzliches Kriterium. Während klassische Batch-ETL-Prozesse mit Stunden- oder Minutentakten arbeiten (vgl. Pfandzelter \& Bermbach, zit.~n. Freitas et al.~2025, S.~808), müssen Sport-Events innerhalb von Sekunden verarbeitet werden. Dies korrespondiert mit der in Kapitel~\ref{motivation-und-anforderungen} formulierten Qualitätsanforderung N4 (P95 unter 30\,s). Das erfordert Plattformen mit niedrigem Overhead, asynchroner Verarbeitung und effizienten Webhook-Mechanismen (vgl. Nakavisute \& Sincharoonsak 2025). Wie diese Anforderungen im System umgesetzt werden, erläutert Abschnitt~\ref{workflow-konzeption-n8n}.



\section{Webanwendungs- und Infrastruktur-Stack}\label{backend-frontend-und-deployment-technologien}

Der Abschnitt gliedert sich in Backend-Framework und Persistenzschicht (Abschnitt~\ref{backend-framework-und-persistenzschicht}), Echtzeit-Kommunikationsmuster (Abschnitt~\ref{echtzeit-kommunikation-im-web}), Frontend-Framework (Abschnitt~\ref{frontend-framework}) und Cloud-Deployment (Abschnitt~\ref{cloud-deployment}).

\subsection{Backend-Framework und Persistenzschicht}\label{backend-framework-und-persistenzschicht}

Im Python-Backend-Ökosystem sind mit \textbf{Flask}, \textbf{Django} und \textbf{FastAPI} drei etablierte Optionen verfügbar (vgl. Grigorev 2019). Für die vorliegende Anwendung ist dabei vor allem das Asynchronitätsmodell entscheidend: Während Flask und Django historisch primär synchron ausgerichtet sind, unterstützt FastAPI den ASGI-Standard und \texttt{asyncio} nativ (vgl. Grigorev 2019). Auf Basis dieser Kriterien wird für die vorliegende Arbeit FastAPI als Backend-Framework gewählt; die architektonische Ausgestaltung wird in Abschnitt~\ref{backend-konzeption} konkretisiert.

Für die Persistenzschicht bieten sich relationale Datenbanksysteme (z.\,B. \textbf{PostgreSQL}, MySQL) und dokumentenorientierte Stores (z.\,B. \textbf{MongoDB}) als Hauptklassen an. Relationale Systeme erzwingen über Fremdschlüsselbeziehungen referenzielle Integrität und unterstützen Joins über normalisierte Schemata. Dokumentenorientierte Stores bieten höhere Schemaflexibilität auf Kosten von ACID-Garantien und deklarativer Join-Semantik (vgl. Fowler \& Sadalage 2012, S.~26--27).

\subsection{Echtzeit-Kommunikationsmuster}\label{echtzeit-kommunikation-im-web}

Für Echtzeitanwendungen wie Liveticker existieren drei etablierte Kommunikationsmuster. \textbf{HTTP-Polling} ist einfach und kompatibel, erzeugt aber inhärente Latenz und Overhead bei leeren Antworten (vgl. de la Torre et al.~2019, S.~261). \textbf{Server-Sent Events (SSE)} ermöglichen einen HTTP-basierten unidirektionalen Push vom Server ohne Protokollwechsel und sind ressourceneffizienter als Polling (vgl. ebd., S.~261--262). \textbf{WebSocket} bietet eine bidirektionale, persistente Verbindung mit minimalem Overhead nach initialem Handshake, ist latenzoptimal, aber ressourcenintensiver (vgl. RFC 6455, de la Torre et al.~2019, S.~262).

Die Muster unterscheiden sich damit primär in Latenz, Ressourcenverbrauch und Übertragungsrichtung. Hybridansätze kombinieren mehrere Muster für unterschiedliche Datenströme eines Systems. Die im Projekt gewählte Kommunikationsarchitektur und ihre Begründung werden in Abschnitt~\ref{kommunikations--und-triggerarchitektur} ausgeführt.

\subsection{Frontend-Framework}\label{frontend-framework}

Für das Frontend wird \textbf{React} als komponentenbasiertes Single-Page-Framework eingesetzt. Ausschlaggebend sind die gute Eignung für asynchrone Live-Datenpfade, die klare Komponentenzerlegung für redaktionelle Oberflächen und die breite Verfügbarkeit etablierter Test- und Build-Werkzeuge im Node-Ökosystem.

\textbf{TypeScript} wird ergänzend als Typschicht eingesetzt. Es reduziert Schnittstellenfehler zwischen Komponenten, API-Responses und Zustandsobjekten bereits zur Kompilierzeit statt erst zur Laufzeit (vgl. Microsoft 2023) und unterstützt damit die Robustheit einer ereignisgetriebenen Live-Oberfläche.

\subsection{Cloud-Deployment}\label{cloud-deployment}

\textbf{Platform-as-a-Service (PaaS)} abstrahiert die Infrastrukturverwaltung und ermöglicht Deployments direkt aus Git-Branches. Es setzt jedoch ein \textbf{Stateless-Design} voraus: Persistente Daten werden in externe Dienste ausgelagert, sodass zusätzliche Instanzen ohne Zustandskonsistenzprobleme gestartet werden können. Als Deployment-Einheit hat sich \textbf{Docker} als Standard etabliert (vgl. Merkel 2014): Container kapseln Anwendung und Laufzeitumgebung in einem reproduzierbaren Image, das plattformunabhängig lauffähig ist und von PaaS-Anbietern direkt aus einem Dockerfile bereitgestellt werden kann. Die konkrete Plattformwahl ist in Abschnitt~\ref{deployment-konzept} begründet.

Zusammenfassend liefert Kapitel~\ref{stand-der-technik} damit die technologiebezogenen Entscheidungsgrundlagen für die Systemkonzeption: LLM-basierte Generierung mit Prompt-Steuerung und HITL-Absicherung, workflowbasierte Echtzeitintegration über ETL-Orchestrierung sowie eine webbasierte, cloudfähige Systemarchitektur. Auf dieser Grundlage konkretisiert Kapitel~\ref{systemkonzeption} die Architekturentscheidungen entlang der in Kapitel~\ref{motivation-und-anforderungen} abgeleiteten Anforderungen.

